\section{Evaluation Methodology}
\label{sec:eval}

\subsection{Attack success rate}

For $n$ independent attack attempts of which $k$ achieve the attacker's goal, the \emph{attack success rate} is $\widehat{\mathrm{ASR}} = k/n$. It is specialized per task: suppression rate (classification), omission rate (summarization), and unsafe-recommendation rate (remediation).

\subsection{Utility and the robustness--utility trade-off}

A defense that makes the copilot label everything as malicious has zero ASR and no value. Every defense must therefore also be measured by its \emph{utility}: task performance on a clean, labelled, un-attacked dataset~\cite{debenedetti2024agentdojo}. For classification, utility is the accuracy or the $F_1$ score of the malicious class (the latter is preferable because benign events usually dominate); for summarization, the fraction of summaries that mention every relevant event; for remediation, the fraction of recommendations judged appropriate. A defense is only useful if it lowers ASR while keeping utility close to that of the undefended model.

\subsection{Confidence intervals for proportions}

$\widehat{\mathrm{ASR}}$ is an estimate of an unknown proportion and must be reported with its uncertainty. The Clopper--Pearson interval~\cite{clopper1934} is the exact interval for a binomial proportion. At level $1-\alpha$ its bounds are quantiles of the Beta distribution:
\begin{equation}
\begin{split}
p_{\mathrm{low}} &= B\!\left(\tfrac{\alpha}{2};\; k,\; n-k+1\right),\\
p_{\mathrm{high}} &= B\!\left(1-\tfrac{\alpha}{2};\; k+1,\; n-k\right),
\end{split}
\label{eq:cp}
\end{equation}
where $B(q; a, b)$ is the $q$-quantile of $\mathrm{Beta}(a,b)$ (with $p_{\mathrm{low}}=0$ if $k=0$ and $p_{\mathrm{high}}=1$ if $k=n$).

For example, with $k=6$ successes out of $n=20$ attempts, $\widehat{\mathrm{ASR}}=0.30$, but the 95\% interval is approximately $[0.12, 0.54]$: the true rate could be anywhere from rare to more than half. With $n=200$ and $k=60$ the point estimate is the same, but the interval shrinks to approximately $[0.24, 0.37]$. In practice, two configurations (e.g.\ with and without a defense) are considered different only if their intervals do not overlap; overlapping intervals are reported as inconclusive rather than as an improvement.

\subsection{Staged experimental design}

The full experimental grid (models $\times$ attack classes $\times$ defenses $\times$ tasks $\times$ fields) contains hundreds of cells, and running each with a large $n$ is too expensive on local hardware. The experiments are therefore run in two stages. In the \emph{screening} stage every cell is run with a small $n$, which is enough to separate clearly ineffective configurations from promising ones. In the \emph{confirmation} stage only the interesting cells (those with high ASR or overlapping intervals between competing configurations) are re-run with a large $n$, so that the final conclusions rest on narrow confidence intervals.
